\chapter{\texorpdfstring{$C^*$}{C*}-algebras: spectral rigidity and functional calculus}

\section{The \texorpdfstring{$C^*$}{C*}-identity and special classes of elements}

\begin{cthdefinicion}{($C^*$-algebra)}
A $C^*$-algebra is an involutive Banach algebra $\A$ satisfying
\[
 \|a^*a\|=\|a\|^2,
 \qquad a\in\A.
\]
\end{cthdefinicion}
\origen{Definition 2.1}

\begin{cthejemplo}{(Operators on a Hilbert space and the $C^*$-identity)}
We seek a concrete model of the definition. If $H$ is a
Hilbert space, $\mathcal B(H)$ is the algebra of bounded linear operators
$H\to H$, with involution given by the adjoint $T\mapsto T^*$.
For $T\in\mathcal B(H)$, the inequality
$\|T^*T\|\le\|T^*\|\|T\|=\|T\|^2$ and the identity
\[
 \|T\|^2=\sup_{\|x\|=1}\langle T^*Tx,x\rangle
 \le\|T^*T\|
\]
give $\|T^*T\|=\|T\|^2$. Thus the operator norm and the adjoint
satisfy precisely the $C^*$-identity.
\end{cthejemplo}

\begin{cthobservacion}{(The identity is compatible with the involution)}
If a $C^*$-algebra $\A$ has an identity $1_\A$, then
$1_\A^*=1_\A$ and $\|1_\A\|=1$.
\end{cthobservacion}
\origen{Observation 2.2}

\begin{proof}[Proof]
For every $a\in\A$, the identity $(a1_\A)^*=1_\A^*a^*=a^*$ shows that
$1_\A^*$ acts as a left identity. Likewise, it acts as a right identity;
by uniqueness of the identity, $1_\A^*=1_\A$. The $C^*$-identity gives
$\|1_\A\|^2=\|1_\A^*1_\A\|=\|1_\A\|$, and $1_\A\ne0$ implies
$\|1_\A\|=1$.
\end{proof}

\begin{cthejemplo}{(A self-adjoint identity of norm one in matrix algebras)}
Let us illustrate the preceding observation with the identity of
$M_n(\mathbb C)$, $n\ge1$, equipped with the operator norm and the
conjugate-transpose adjoint. The matrix $I_n$ satisfies
$I_n^*=I_n$; moreover, $\|I_nx\|=\|x\|$ for every vector, so
$\|I_n\|=1$. The abstract assertion about $1_\A$ becomes here
a visible property of the identity matrix.
\end{cthejemplo}

\begin{investigacion}
\textbf{Direct example.} For the open logarithmic star $\Gamma$ in the Haseman
problem, $C_b(\Gamma)$, with the uniform norm and pointwise conjugation, is a
unital $C^*$-algebra; the subalgebra $\SO(\Gamma)$ consists of functions that
oscillate slowly at $0$ and $\infty$ \cite[§1, p.~464]{Karlovich2017}.
\textbf{What to recognize:} both are function algebras of the type just
defined; this observation only illustrates the $C^*$-identity and plays no
part in the course proofs.
\end{investigacion}

\begin{cthdefinicion}{($*$-homomorphism)}
Let $\A,\B$ be $C^*$-algebras. A $*$-homomorphism is a linear, multiplicative
homomorphism $\Phi:\A\to\B$ satisfying
\[
 \Phi(a^*)=\Phi(a)^*\qquad(a\in\A).
\]
If both algebras are unital, we say that $\Phi$ is unital if $\Phi(1_\A)=1_\B$.
\end{cthdefinicion}
\origen{Definition 2.4}

\begin{cthejemplo}{(Functions as multiplication operators)}
Let $X$ be locally compact and Hausdorff. We want to represent the algebra
$C_b(X)$ by operators on $\ell^2(X)$, the space of square-summable complex
families indexed by $X$. For $f\in C_b(X)$ we define
$(M_f\xi)(x)=f(x)\xi(x)$. The estimate
$\|M_f\xi\|_2\le\|f\|_\infty\|\xi\|_2$ proves that $M_f$ is bounded.
Moreover, $M_fM_g=M_{fg}$ and, by comparing inner products,
$M_f^*=M_{\overline f}$; also $M_1=I$.
Hence $f\mapsto M_f$ is a unital $*$-homomorphism.
Problem 2.1 checks that it even preserves the norm: a function algebra
can be represented faithfully as an operator algebra.
\end{cthejemplo}

\begin{cthdefinicion}{(Normal, self-adjoint, unitary, and projection)}
Let $\A$ be a $C^*$-algebra and $a\in\A$.
\begin{enumerate}[label=(\roman*)]
\item $a$ is \emph{normal} if $a^*a=aa^*$;
\item $a$ is \emph{self-adjoint} if $a^*=a$; we write
$\A_{\mathrm{sa}}=\{x\in\A:x=x^*\}$;
\item if $\A$ has an identity $1_\A$, $a$ is \emph{unitary} if
$a^*a=aa^*=1_\A$; we write $\mathcal U(\A)$ for the set of unitaries;
\item $a$ is a \emph{projection} if $a=a^*=a^2$; we write
$\mathcal P(\A)$ for the set of projections.
\end{enumerate}
\end{cthdefinicion}
\origen{Definition 2.5}

\begin{cthejemplo}{(Normality, self-adjointness, unitarity, and projection in $M_2$)}
Let us find a representative of each class in the definition in
$M_2(\mathbb C)$:
\[
\begin{array}{c|c}
\text{type}&\text{example}\\ \hline
\text{normal}&\begin{pmatrix}1&0\\0&i\end{pmatrix}\\[2mm]
\text{self-adjoint}&\begin{pmatrix}1&0\\0&-2\end{pmatrix}\\[2mm]
\text{unitary}&\begin{pmatrix}0&1\\1&0\end{pmatrix}\\[2mm]
\text{projection}&\begin{pmatrix}1&0\\0&0\end{pmatrix}
\end{array}
\]
The two diagonal matrices commute with their adjoints; the second
equals its adjoint. For the third, $U^*=U$ and $U^2=I_2$,
so $U^*U=UU^*=I_2$. The fourth satisfies
$P=P^*=P^2$. These checks show why every
self-adjoint element, unitary, or projection is normal, although the first
normal matrix is not self-adjoint.
\end{cthejemplo}

\begin{investigacion}
\textbf{Direct example in $L^2$; an interpretive analogy for Haseman.}
For $\Gamma_0=\mathbb R$ or the circle
$\mathbb T=\{z\in\mathbb C:|z|=1\}$, with arc-length measure on $L^2(\Gamma_0)$,
the Cauchy operator satisfies $S_{\Gamma_0}^2=I$ and defines
\[
 P_{\Gamma_0}^\pm=\frac12(I\pm S_{\Gamma_0}).
\]
The identity $S_{\Gamma_0}^2=I$ gives $(P_{\Gamma_0}^\pm)^2=P_{\Gamma_0}^\pm$ by
algebraic calculation. In this model, moreover, $S_{\Gamma_0}$ is self-adjoint and
$(P_{\Gamma_0}^\pm)^*=P_{\Gamma_0}^\pm$: these are orthogonal projections
\cite[Sec.~2, p.~22]{BFK2006}. The formula
$P_\pm=(I\pm S_\Gamma)/2$ reappears in the Haseman operator on
$L^p(\Gamma)$ for an open logarithmic star $\Gamma$, with
$1<p<\infty$ and data satisfying the hypotheses of \cite{Karlovich2017},
\[
 T=V_\alpha P_+ + GP_-\in\mathcal B(L^p(\Gamma)),
\]
which applies the shift $V_\alpha f=f\circ\alpha$ to the $P_+f$ term and
multiplication by $G$ to the $P_-f$ term. On the open star,
$S_\Gamma^2=I$ and orthogonality do not follow from the $\Gamma_0$ case
\cite[(1.2)]{Karlovich2017}.
\textbf{What to recognize:} $S^2=I$ ensures idempotence in the Hilbert-space
model and $S^*=S$ adds orthogonality; the Haseman formula is read term by
term without transferring these properties to the open star.
\end{investigacion}

\pagebreak
\section{Norm and spectrum of normal elements}

\begin{cthteorema}{(The norm of a normal element equals its spectral radius)}
Let $\A$ be a unital $C^*$-algebra and $a\in\A$ normal. Then
\[
 \|a\|=r(a).
\]
\end{cthteorema}
\origen{Theorem 2.6}

\begin{proof}[Proof]
Normality allows us to rearrange the factors:
\[
\|a^2\|^2
=\|(a^2)^*a^2\|
=\|(a^*)^2a^2\|
=\|(a^*a)^2\|
=\|a^*a\|^2
=\|a\|^4.
\]
Thus $\|a^2\|=\|a\|^2$. Since $a^2$ is also normal, induction gives
$\|a^{2^n}\|=\|a\|^{2^n}$. Applying the spectral radius formula along the
subsequence $2^n$,
\[
 r(a)=\lim_{n\to\infty}\|a^{2^n}\|^{1/2^n}=\|a\|.
\]
\end{proof}

\proofdiagram
{$\A$ is unital $C^*$; $a\in\A$ normal}
{$C^*$-identity; spectral radius formula}
{Prove $\|a^{2^n}\|=\|a\|^{2^n}$ by induction}
{$r(a)=\lim_n\|a^{2^n}\|^{1/2^n}=\|a\|$}

\begin{cthejemplo}{(Norm and radius of a normal matrix)}
Let us check the theorem for $A=\operatorname{diag}(2,-i)$ with the
Euclidean operator norm. Since $A$ is diagonal, it commutes with
$A^*$ and is normal. Its spectral values are $2$ and $-i$;
therefore $r(A)=\max\{2,1\}=2$. The norm of a diagonal matrix
in this norm is the largest absolute value of its entries, so
$\|A\|=2=r(A)$. This calculation makes concrete the equality that the
theorem guarantees for every normal element.
\end{cthejemplo}

\begin{cthnota}{(Where the hypotheses are used)}
Normality allows us to pass from $(a^*)^2a^2$ to $(a^*a)^2$.
The $C^*$-identity turns that algebraic equality into an equality of norms;
the spectral radius formula, proved in Part 1 for Banach algebras,
completes the argument.
\end{cthnota}


\begin{cthcorolario}{(The $C^*$-norm is determined algebraically)}
Let $\A$ be a unital $C^*$-algebra. For every $a\in\A$,
\[
 \|a\|=r(a^*a)^{1/2}.
\]
In particular, two norms that make the same unital involutive algebra a
$C^*$-algebra coincide.
\end{cthcorolario}
\origen{Corollary 2.7}

\begin{proof}[Proof]
The element $a^*a$ is self-adjoint and hence normal. Therefore
$r(a^*a)=\|a^*a\|=\|a\|^2$ by the preceding theorem and the $C^*$-identity.
Taking nonnegative square roots gives the stated formula.
If $\|\cdot\|_1$ and $\|\cdot\|_2$ make the same unital involutive algebra a
$C^*$-algebra, the invertibility of $a^*a-\lambda1_\A$ depends only on the
product and the identity. Thus $r(a^*a)$ is the same for both norms, and the
formula gives $\|a\|_1=\|a\|_2$ for every $a\in\A$.
\end{proof}

\proofdiagram
{$\A$ is unital $C^*$; $a\in\A$; $a^*a$ is normal}
{Previous theorem + $C^*$-identity}
{$r(a^*a)=\|a^*a\|=\|a\|^2$}
{$\|a\|=r(a^*a)^{1/2}$; uniqueness of the norm}

\begin{cthejemplo}{(The uniform norm is forced on $C(X)$)}
We want to apply the norm-uniqueness corollary to a function algebra.
Let $X$ be compact Hausdorff and $f\in C(X)$, with pointwise conjugation
as involution. The element $f^*f=|f|^2$ has spectrum $|f|^2(X)$,
by the criterion of pointwise invertibility from Part 1. The corollary gives
\[
 \|f\|=r(f^*f)^{1/2}
 =\bigl(\max_{x\in X}|f(x)|^2\bigr)^{1/2}
 =\|f\|_\infty.
\]
Any $C^*$-norm compatible with these operations must therefore
coincide with the uniform norm.
\end{cthejemplo}

\begin{cthteorema}{(Unital $*$-homomorphisms are contractive)}
Let $\A,\B$ be unital $C^*$-algebras and $\Phi:\A\to\B$ a unital $*$-homomorphism. Then
\[
 \|\Phi(a)\|\le\|a\|\quad(a\in\A),
 \qquad \|\Phi\|=1.
\]
\end{cthteorema}
\origen{Theorem 2.8}

\begin{proof}[Proof]
For each $b\in\A$, the unital homomorphism satisfies
$\spec(\Phi(b))\subset\spec(b)$ and hence
$r(\Phi(b))\le r(b)$. Taking $b=a^*a$,
\[
\|\Phi(a)\|^2
=\|\Phi(a)^*\Phi(a)\|
=\|\Phi(a^*a)\|
=r(\Phi(a^*a))
\le r(a^*a)=\|a\|^2.
\]
Since $\Phi(1_\A)=1_\B$ and both identities have norm one,
$\|\Phi\|\ge1$; contractivity gives equality.
\end{proof}

\proofdiagram
{$\Phi$ is unital and a $*$-homomorphism}
{Spectral inclusion for homomorphisms; normal-element theorem $\|x\|=r(x)$}
{Apply to $a^*a$: $\Phi(a^*a)=\Phi(a)^*\Phi(a)$}
{$\|\Phi(a)\|^2\le\|a\|^2$ and $\|\Phi\|=1$}

\begin{cthejemplo}{(Restriction as a contractive $*$-homomorphism)}
Let $X$ be compact Hausdorff and $Y\subset X$ closed and nonempty.
We want to compare a function defined on $X$ with its restriction to
$Y$. The map $\Phi:C(X)\to C(Y)$,
$\Phi(f)=f|_Y$, preserves sums, products, pointwise conjugation, and the
constant function $1$, so it is a unital $*$-homomorphism.
The contractivity theorem implies
$\|f|_Y\|_\infty\le\|f\|_\infty$, an inequality also obtained by
taking the supremum over a smaller set. Restricting the domain cannot
increase the uniform norm.
\end{cthejemplo}

\begin{cthobservacion}{(Real and imaginary parts)}
Let $\A$ be a $C^*$-algebra. Every $a\in\A$ has a unique expression
\[
 a=x+iy,
 \qquad
 x=\frac{a+a^*}{2},\quad y=\frac{a-a^*}{2i},
\]
with $x=x^*$ and $y=y^*$.
\end{cthobservacion}
\origen{Observation 2.9}

\begin{proof}[Proof]
The given formulas satisfy $x^*=x$, $y^*=y$, and $x+iy=a$. If also
$a=u+iv$ with $u=u^*$ and $v=v^*$, then $a^*=u-iv$; adding and subtracting
these two identities yields $u=(a+a^*)/2=x$ and
$v=(a-a^*)/(2i)=y$. This proves uniqueness.
\end{proof}

\begin{cthejemplo}{(Self-adjoint parts of a non-Hermitian matrix)}
To see the decomposition in the observation in a nontrivial case,
take $A=\begin{pmatrix}1&i\\0&2\end{pmatrix}$.
Calculating the real and imaginary parts gives
\[
 X=\frac{A+A^*}{2}
 =\begin{pmatrix}1&i/2\\-i/2&2\end{pmatrix},
 \qquad
 Y=\frac{A-A^*}{2i}
 =\begin{pmatrix}0&1/2\\1/2&0\end{pmatrix}.
\]
We check that $X^*=X$, $Y^*=Y$, and $A=X+iY$.
Thus even a matrix that is not self-adjoint can be reconstructed from
two matrices that are.
\end{cthejemplo}

\begin{cthlema}{(Vertical estimate for self-adjoint elements)}
Let $\A$ be a unital $C^*$-algebra, $x\in\A$ self-adjoint, and $t\in\mathbb R$. Then
\[
 \|x+it1_\A\|^2\le \|x\|^2+t^2.
\]
\end{cthlema}
\origen{Lemma 2.10}

\begin{proof}[Proof]
\[
\|x+it1_\A\|^2
=\|(x-it1_\A)(x+it1_\A)\|
=\|x^2+t^2 1_\A\|
\le\|x^2\|+t^2
=\|x\|^2+t^2.
\]
\end{proof}

\proofdiagram
{$x\in\A$ self-adjoint; $t\in\mathbb R$}
{$C^*$-identity}
{$\|x+it1_\A\|^2=\|x^2+t^2 1_\A\|$}
{$\le\|x\|^2+t^2$}

\begin{cthejemplo}{(The vertical bound is sharp for real scalars)}
Let us check whether the lemma's inequality can be improved in general.
In the $C^*$-algebra $\mathbb C$, take the self-adjoint element
$s\in\mathbb R$ and $t\in\mathbb R$. Then
$|s+it|^2=(s-it)(s+it)=s^2+t^2$ and $\|s\|^2=s^2$.
The lemma's inequality is an equality in this case; the term
$t^2$ and the coefficient of $\|s\|^2$ are therefore natural.
\end{cthejemplo}

\begin{cthproposicion}{(A vertical bound forces a scalar to be real)}
If $z\in\mathbb C$ and there exists $C\ge0$ such that
\[
 |z+it|^2\le C+t^2\qquad(t\in\mathbb R),
\]
then $z\in\mathbb R$.
\end{cthproposicion}
\origen{Proposition 2.11}

\begin{proof}[Proof]
If $z=x+iy$ with $x,y\in\mathbb R$, then
$|z+it|^2-t^2=x^2+y^2+2yt$. If $y>0$, the right-hand side tends to $+\infty$
as $t\to+\infty$; if $y<0$, this happens as $t\to-\infty$. The existence
of a uniform bound forces $y=0$.
\end{proof}

\proofdiagram
{$z\in\mathbb C$; $C\ge0$; write $z=x+iy$}
{Expand $|z+it|^2-t^2=x^2+y^2+2yt$}
{If $y\neq0$, let $t\to\pm\infty$ with the sign of $y$}
{The bound fails; hence $y=0$}

\begin{cthejemplo}{(A bound that fails for $1+i$)}
To see how Proposition~2.11 detects an imaginary part, test its hypothesis with
$z=1+i$. The calculation gives
\[
 |z+it|^2-t^2=|1+i(1+t)|^2-t^2=2+2t.
\]
As $t\to+\infty$, this expression exceeds every constant $C$. Thus there is
no bound $|1+i+it|^2\le C+t^2$ valid for every $t\in\mathbb R$.
The term linear in $t$ is the concrete sign that $z$ is not real.
\end{cthejemplo}

\begin{cthproposicion}{(Functionals attaining their norm at the identity are real on self-adjoint elements)}
Let $\A$ be a unital $C^*$-algebra and $\varphi:\A\to\mathbb C$ a continuous
linear functional with $\|\varphi\|=\varphi(1_\A)$. Then
\[
 \varphi(x)\in\mathbb R\qquad\text{for every }x\in\A\text{ with }x=x^*.
\]
\end{cthproposicion}
\origen{Proposition 2.12}

\begin{proof}[Proof]
If $\varphi=0$, then $\varphi(x)=0\in\mathbb R$ for every self-adjoint $x$.
Otherwise, $\varphi(1_\A)=\|\varphi\|>0$.
Define $\psi=\varphi/\varphi(1_\A)$; then $\|\psi\|=\psi(1_\A)=1$.
For $t\in\mathbb R$,
\[
 |\psi(x)+it|^2
 =|\psi(x+it1_\A)|^2
 \le\|x+it1_\A\|^2
 \le\|x\|^2+t^2.
\]
The previous proposition, with $z=\psi(x)$, implies that $\psi(x)$ is real;
since $\varphi(1_\A)>0$, so is $\varphi(x)$.
\end{proof}

\proofdiagram
{$\varphi:\A\to\mathbb C$ continuous; $\|\varphi\|=\varphi(1_\A)$; $x=x^*$}
{Vertical estimate lemma + preceding scalar proposition}
{If $\varphi\ne0$: $\psi=\varphi/\varphi(1_\A)$ and $|\psi(x)+it|^2\le\|x+it1_\A\|^2$}
{$\varphi=0$ or $\psi(x)\in\mathbb R$; hence $\varphi(x)\in\mathbb R$}

\begin{cthejemplo}{(An evaluation satisfies the reality hypothesis)}
Take $X=[0,1]$ and $x_0\in X$. The evaluation
$\varphi=\operatorname{ev}_{x_0}:C(X)\to\mathbb C$ satisfies
$\|\varphi\|=\varphi(1)=1$, as checked in Section~1.
For the self-adjoint function $h(t)=t^2$, we obtain
$\varphi(h)=x_0^2\in\mathbb R$.
Proposition~2.12 gives the same conclusion for \emph{every} self-adjoint
element of $C(X)$ and, more generally, for functionals with this norm
condition in any unital $C^*$-algebra.
\end{cthejemplo}

\begin{cthproposicion}{(The spectrum of a self-adjoint element is real)}
Let $\A$ be a unital $C^*$-algebra and $a\in\A$ self-adjoint. Then
\[
 \spec(a)\subset\mathbb R.
\]
\end{cthproposicion}
\origen{Proposition 2.13}

\begin{proof}[Proof]
By Problem 1.9 there is a commutative, involutive, unital Banach subalgebra
$\mathcal C\subset\A$ that contains $a$ and preserves its spectrum. In it,
\[
 \spec_\A(a)=\spec_{\mathcal C}(a)=\{\varphi(a):\varphi\in\Omega(\mathcal C)\}.
\]
Since $\mathcal C$ is closed and inherits the $C^*$-identity, it too is a
$C^*$-algebra. Every character of $\mathcal C$ has norm one and
$\varphi(1_{\mathcal C})=1$, so the previous proposition gives
$\varphi(a)\in\mathbb R$. The spectral description above then shows
that every $\lambda\in\spec_\A(a)$ is real.
\end{proof}

\proofdiagrambranch
{$\A$ is unital $C^*$; $a\in\A$, $a=a^*$}
{Reduce to a commutative subalgebra preserving the spectrum (Problem 1.9)}
{Spectrum via characters + functionals real on self-adjoint elements}
{$\lambda\in\spec(a)\Rightarrow\lambda=\varphi(a)\in\mathbb R$}
{$\spec(a)\subset\mathbb R$}

\begin{cthejemplo}{(Real eigenvalues of a Hermitian matrix)}
Consider $H=\begin{pmatrix}2&i\\-i&3\end{pmatrix}\in M_2(\mathbb C)$.
Conjugate transposition gives $H^*=H$, so Proposition~2.13 predicts
$\spec(H)\subset\mathbb R$. Indeed,
\[
 \det(\lambda I-H)=(\lambda-2)(\lambda-3)-1
 =\lambda^2-5\lambda+5,
 \qquad
 \spec(H)=\left\{\frac{5-\sqrt5}{2},\frac{5+\sqrt5}{2}\right\}.
\]
This finite-dimensional calculation confirms the spectral conclusion
that the proposition extends to self-adjoint elements with no associated matrix.
\end{cthejemplo}

\begin{cthnota}{(Where the hypotheses are used)}
The equality $a=a^*$ allows us to apply Problem 1.9 in its involutive case
and then the proposition on real-valued functionals. The identity is used
in the description of the spectrum by normalized characters.
\end{cthnota}


\begin{cthdefinicion}{(Positive element and order)}
Let $\A$ be a unital $C^*$-algebra. An element $x\in\A$ is \emph{positive},
written $x\ge0$, if $x=x^*$ and $\spec_\A(x)\subset[0,\infty)$.
For self-adjoint $x,y\in\A$, $x\le y$ means that $y-x\ge0$.
\end{cthdefinicion}

\begin{puente}
This is one of the clearest connections between the problems and Section 2:
the original proof that the spectrum is real explicitly uses the exercise in
Problem 1 that constructs a commutative subalgebra containing the element and
preserving its spectrum. The exercise is an essential part of the proof.
\end{puente}

\section{Characters, spectral permanence, and representation}

\begin{cthcorolario}{(Functionals attaining their norm at the identity preserve the involution)}
Let $\A$ be a unital $C^*$-algebra and $\varphi:\A\to\mathbb C$ a continuous
linear functional with $\|\varphi\|=\varphi(1_\A)$. Then, for every $a\in\A$,
\[
 \varphi(a^*)=\overline{\varphi(a)}.
\]
In particular, this holds for the characters of $\A$.
\end{cthcorolario}
\origen{Corollary 2.14}

\begin{proof}[Proof]
Let $a\in\A$ and write $a=x+iy$, where
$x=(a+a^*)/2$ and $y=(a-a^*)/(2i)$ are self-adjoint. By the previous
proposition, $\varphi(x),\varphi(y)\in\mathbb R$, and
\[
 \varphi(a^*)=\varphi(x-iy)=\varphi(x)-i\varphi(y)
 =\overline{\varphi(x)+i\varphi(y)}=\overline{\varphi(a)}.
\]
If $\varphi$ is a character, then it is a continuous functional with
$\varphi(1_\A)=\|\varphi\|=1$, so it satisfies the hypothesis used.
\end{proof}

\proofdiagram
{$\varphi:\A\to\mathbb C$ continuous; $\|\varphi\|=\varphi(1_\A)$; $a=x+iy$ with $x,y$ self-adjoint}
{Proposition 2.12: $\varphi(x),\varphi(y)\in\mathbb R$}
{$\varphi(a^*)=\varphi(x)-i\varphi(y)$}
{$\varphi(a^*)=\overline{\varphi(a)}$}

\begin{cthejemplo}{(An evaluation preserves conjugation)}[unbreakable]
In $C([0,1])$ choose $f(t)=1+it$ and $x\in[0,1]$.
The involution is pointwise conjugation, so
$f^*(t)=1-it$. For the character $\operatorname{ev}_x$,
\[
 \operatorname{ev}_x(f^*)=1-ix
 =\overline{1+ix}
 =\overline{\operatorname{ev}_x(f)}.
\]
This equality illustrates Corollary~2.14 for a concrete character; the
corollary also guarantees it for functionals attaining their norm at the identity.
\end{cthejemplo}

\begin{cthcorolario}{(The Gelfand transform is a $*$-homomorphism)}
If $\A$ is a unital commutative $C^*$-algebra, then
\[
 \Gamma:\A\to C(\Omega(\A)),\qquad a\mapsto\widehat a,
\]
is a unital $*$-homomorphism.
\end{cthcorolario}
\origen{Corollary 2.15}

\begin{proof}[Proof]
For $a,b\in\A$ and $\varphi\in\Omega(\A)$,
$\widehat{a+b}(\varphi)=\widehat a(\varphi)+\widehat b(\varphi)$,
$\widehat{ab}(\varphi)=\widehat a(\varphi)\widehat b(\varphi)$, and
$\widehat{1_\A}(\varphi)=1$, by linearity, multiplicativity, and the unital
property of the character. Also
$\widehat{\lambda a}(\varphi)=\lambda\widehat a(\varphi)$ for
$\lambda\in\mathbb C$. Thus the Gelfand transform is a linear,
multiplicative, unital homomorphism. Every $\varphi\in\Omega(\A)$ satisfies
$\|\varphi\|=\varphi(1_\A)=1$, so the previous corollary gives, for $a\in\A$,
\[
 \widehat{a^*}(\varphi)=\varphi(a^*)
 =\overline{\varphi(a)}=(\widehat a)^*(\varphi).
\]
Since the equality holds for every $\varphi\in\Omega(\A)$, the transform
preserves $*$.
\end{proof}

\proofdiagram
{Gelfand is a unital homomorphism}
{Characters preserve $*$}
{$\widehat{a^*}(\varphi)=\varphi(a^*)=\overline{\varphi(a)}$}
{$\widehat{a^*}=(\widehat a)^*$}

\begin{cthejemplo}{(Conjugation under the Gelfand transform)}
Take $\A=C([0,1])$ and $f(t)=1+it$. Since the characters are evaluations,
$\widehat f(\operatorname{ev}_x)=1+ix$ for each $x\in[0,1]$. Likewise,
\[
 \widehat{f^*}(\operatorname{ev}_x)=1-ix
 =\overline{\widehat f(\operatorname{ev}_x)}.
\]
Corollary~2.15 states this identity simultaneously for all elements and
characters: the Gelfand transform respects the involution.
\end{cthejemplo}

\begin{cthproposicion}{(Invertibility detected by $a^*a$ and $aa^*$)}
For $a$ in a unital $C^*$-algebra $\A$,
\[
 a\in\GL(\A)
 \quad\Longleftrightarrow\quad
 a^*a\in\GL(\A)\text{ and }aa^*\in\GL(\A).
\]
\end{cthproposicion}
\origen{Proposition 2.16}

\begin{proof}[Proof]
If $a\in\GL(\A)$, the identity
$(a^{-1})^*a^*=a^*(a^{-1})^*=1_\A$ shows that $a^*$ is invertible.
Hence its products $a^*a$ and $aa^*$ are invertible. Conversely, take
$y=(a^*a)^{-1}$ and $z=(aa^*)^{-1}$, so that
$y(a^*a)=1_\A$ and $(aa^*)z=1_\A$. Then $ya^*$ is a left inverse of $a$ and
$a^*z$ is a right inverse. Moreover,
\[
 ya^*=ya^*(aa^*z)=(ya^*a)a^*z=a^*z.
\]
Therefore a two-sided inverse exists.
\end{proof}

\proofdiagrambranch
{$a\in\A$; $a^*a\in\GL(\A)$}
{$aa^*\in\GL(\A)$}
{$y=(a^*a)^{-1}$ and $z=(aa^*)^{-1}$}
{$ya^*$ and $a^*z$ coincide by associativity}
{$a\in\GL(\A)$}

\begin{cthejemplo}{(A matrix's inverses detected by its products)}
In $M_2(\mathbb C)$, let
$A=\begin{pmatrix}1&1\\0&1\end{pmatrix}$. We have
\[
 A^*A=\begin{pmatrix}1&1\\1&2\end{pmatrix},
 \qquad
 AA^*=\begin{pmatrix}2&1\\1&1\end{pmatrix},
\]
and both determinants are $1$. By Proposition~2.16, $A$ is invertible;
directly, $A^{-1}=\begin{pmatrix}1&-1\\0&1\end{pmatrix}$.
The check with square matrices shows how the two products provide left
and right inverses in the abstract argument.
\end{cthejemplo}

\begin{cthproposicion}{(Spectral permanence in $C^*$-subalgebras)}
Let $\A$ be a unital $C^*$-algebra and $\B\subset\A$ a unital $C^*$-subalgebra
with $1_\B=1_\A$. For every $b\in\B$,
\[
 b\in\GL(\B)\Longleftrightarrow b\in\GL(\A),
 \qquad
 \spec_\B(b)=\spec_\A(b).
\]
\end{cthproposicion}
\origen{Proposition 2.17}

\begin{proof}[Proof]
If $b\in\GL(\B)$, its inverse $b^{-1}\in\B\subset\A$ satisfies
$bb^{-1}=b^{-1}b=1_\B=1_\A$; hence $b\in\GL(\A)$.
For the other direction, first let $h\in\B$ be self-adjoint and invertible
in $\A$. Since $\spec_\B(h)\subset\mathbb R$, for every $n\ge1$ the element
$h+i n^{-1}1_\B$ is invertible in $\B$. These elements converge to $h$;
inversion is continuous on $\GL(\A)$, so their inverses converge in $\A$
to $h^{-1}$. Since $\B$ is closed, $h^{-1}\in\B$.
For a general $b\in\B\cap\GL(\A)$, apply this argument to $b^*b$ and $bb^*$,
which are invertible in $\A$ when $b$ is, and obtain that both are invertible
in $\B$. The previous proposition then concludes that $b$ is invertible in
$\B$. Finally, for each $\lambda\in\mathbb C$ apply the invertibility
equivalence to $b-\lambda1_\B=b-\lambda1_\A$; it follows that
$\lambda\notin\spec_\B(b)$ if and only if
$\lambda\notin\spec_\A(b)$, that is,
$\spec_\B(b)=\spec_\A(b)$.
\end{proof}

\proofdiagrambranch
{$b\in\B\cap\GL(\A)$}
{$h=b^*b$ or $bb^*$: self-adjoint and invertible in $\A$}
{$\spec_\B(h)\subset\mathbb R$; $h+i n^{-1}1_\B\in\GL(\B)$}
{Continuity of inversion in $\A$ and closure of $\B$: $h^{-1}\in\B$}
{$b^*b,bb^*\in\GL(\B)$; hence $b\in\GL(\B)$}

\begin{cthejemplo}{(Counterexample: the disk algebra does not preserve spectra)}
Problem~1.12 compares the disk algebra $A(\mathbb D)$ with $C(\mathbb T)$.
Take the coordinate function $u(z)=z$ on $\mathbb T$. Since
$|u(z)|=1$ on the circle, $u$ is invertible in $C(\mathbb T)$,
with inverse $\overline u$. However, $\overline u$ does not belong to
$A(\mathbb D)$: a holomorphic inverse of $z$ on the disk would have a pole
at $0$. Hence $0\in\spec_{A(\mathbb D)}(u)$ whereas
$0\notin\spec_{C(\mathbb T)}(u)$.
This closed subalgebra is not stable under pointwise involution and does not
satisfy the $C^*$ hypothesis of Proposition~2.17.
\end{cthejemplo}

\begin{cthnota}{(Where the hypotheses are used)}
Stability of $\B$ under involution allows us to use $b^*b$ and $bb^*$ within it;
the common identity lets us compare invertibility, and closedness of $\B$
lets us retain the limit of the inverses.
\end{cthnota}


\begin{cthproposicion}{(The spectrum of a normal element is the character space it generates)}
Let $\A$ be a unital $C^*$-algebra and $a\in\A$ normal. Denote by
$C^*(a,1_\A)$ the smallest unital $C^*$-subalgebra of $\A$ containing $a$. Then
\[
 \Omega(C^*(a,1_\A))\longrightarrow\spec_\A(a),\qquad
 \varphi\longmapsto\varphi(a),
\]
is a homeomorphism.
\end{cthproposicion}
\origen{Proposition 2.18}

\begin{proof}[Proof]
The $C^*$-algebra $C^*(a,1_\A)$ is commutative because $a$ commutes with
$a^*$. By spectral permanence, the spectrum of $a$ in this subalgebra
is the same as in $\A$. In a unital commutative Banach algebra, every point
of the spectrum of $a$ is $\varphi(a)$ for some character $\varphi$; this gives
surjectivity. The map is continuous because the topology of
$\Omega(C^*(a,1_\A))$ is that of pointwise convergence on the algebra's
elements. If $\varphi,\psi\in\Omega(C^*(a,1_\A))$ and
$\varphi(a)=\psi(a)$, then also
$\varphi(a^*)=\overline{\varphi(a)}=\overline{\psi(a)}=\psi(a^*)$.
Since polynomials in $a$ and $a^*$ are dense in $C^*(a,1_\A)$,
linearity and multiplicativity give equality on those polynomials, and
continuity of the characters extends equality to the whole subalgebra.
Thus $\varphi=\psi$. The maximal ideal space and the spectrum are compact
Hausdorff spaces; a continuous bijection between them has a continuous
inverse and is a homeomorphism.
\end{proof}

\proofdiagram
{$a$ normal}
{$C^*(a,1_\A)$ is commutative; spectral permanence; characters}
{Surjective by the spectrum; injective because $a,a^*$ generate densely}
{Continuous bijection between compact Hausdorff spaces: homeomorphism}

\begin{cthejemplo}{(Interpolation on the spectrum of a diagonal matrix)}
Let $\lambda_1,\ldots,\lambda_n$ be distinct values and
$A=\operatorname{diag}(\lambda_1,\ldots,\lambda_n)$.
Proposition~2.18 identifies the characters of $C^*(A,I_n)$ with these
eigenvalues. Concretely, for arbitrary numbers $c_1,\ldots,c_n$,
Lagrange interpolation produces a polynomial $p$ such that
$p(\lambda_j)=c_j$; then
\[
 p(A)=\operatorname{diag}(c_1,\ldots,c_n).
\]
Thus $C^*(A,I_n)$ is the entire diagonal algebra and is identified with
$C(\{\lambda_1,\ldots,\lambda_n\})$. Each character selects one coordinate.
\end{cthejemplo}

\begin{cthteorema}{(Complex Stone--Weierstrass theorem)}
Let $X$ be compact Hausdorff and $\mathcal D\subset C(X)$ a unital
$*$-subalgebra that separates points. Then $\mathcal D$ is dense in $C(X)$
in the uniform norm.
\end{cthteorema}
\origen{Theorem 2.19. Used as an external tool; its full proof is not reproduced here.}

\begin{cthejemplo}{(Polynomials in the coordinate and its conjugate on a compact set)}
Let $K\subset\mathbb C$ be compact. We want to approximate functions in
$C(K)$ by algebraic expressions in the coordinate $z$. The unital
subalgebra of polynomials in $z$ and $\overline z$ separates points because
the function $z\mapsto z$ already distinguishes them; it is also stable
under conjugation. Stone--Weierstrass Theorem~2.19 therefore says that it
is dense in $C(K)$ in the uniform norm.
On $K=\mathbb T$ this includes trigonometric polynomials, since
$\overline z=z^{-1}$. Problem~2.8 uses precisely this density for the
functional calculus.
\end{cthejemplo}

\begin{cthteorema}{(Commutative Gelfand--Naimark theorem)}
Let $\A$ be a unital commutative $C^*$-algebra. The Gelfand transform
\[
 \Gamma:\A\longrightarrow C(\Omega(\A))
\]
is an isometric $*$-isomorphism.
\end{cthteorema}
\origen{Theorem 2.20}

\begin{proof}[Proof]
For each $a\in\A$, commutativity makes $a$ normal. The description of
the spectrum by characters, together with the norm theorem for normal
elements, gives
\[
 \|\widehat a\|_\infty=r(a)=\|a\|.
\]
Thus $\Gamma$ is isometric and injective. Its image is also closed: a
convergent sequence of images comes, by the isometry, from a Cauchy sequence
in $\A$, which converges by completeness. The image is a unital
$*$-subalgebra of $C(\Omega(\A))$ and separates points: if
$\varphi\neq\psi$, there is $a$ with $\varphi(a)\neq\psi(a)$.
Stone--Weierstrass says that the image is dense. Being closed and dense,
it is all of $C(\Omega(\A))$.
\end{proof}

\proofdiagrambranch
{$\A$ commutative unital $C^*$; $a\in\A$}
{Gelfand is a $*$-homomorphism; $\|a\|=r(a)$; Stone--Weierstrass}
{Isometry: $\|\widehat a\|_\infty=r(a)=\|a\|$}
{The image is closed, unital, $*$-stable, and separates points; hence it is dense}
{$\Gamma$ is surjective and isometric}

\begin{cthejemplo}{(Diagonal matrices as functions on two points)}
Let $\A=\{\operatorname{diag}(\alpha,\beta):\alpha,\beta\in\mathbb C\}$.
Its two characters are evaluations of the entries,
$\varphi_1(\operatorname{diag}(\alpha,\beta))=\alpha$ and
$\varphi_2(\operatorname{diag}(\alpha,\beta))=\beta$.
The Gelfand transform of Theorem~2.20 sends
$\operatorname{diag}(\alpha,\beta)$ to the function
$f:\{\varphi_1,\varphi_2\}\to\mathbb C$ with values $\alpha,\beta$.
The equality $\|\operatorname{diag}(\alpha,\beta)\|=\max\{|\alpha|,|\beta|\}
=\|f\|_\infty$ displays the theorem's isometry in this finite model.
\end{cthejemplo}

\begin{cthcorolario}{(Classification of unital commutative $C^*$-algebras)}
Let $\A$ be a unital commutative $C^*$-algebra. There exists a compact
Hausdorff space $X$ such that $\A$ is isometrically $*$-isomorphic to $C(X)$;
this space $X$ is unique up to homeomorphism.
\end{cthcorolario}
\origen{Corollary 2.21}

\begin{proof}[Proof]
Existence follows from Gelfand--Naimark with $X=\Omega(\A)$.
For uniqueness, let $Y$ be compact Hausdorff and
$\varphi\in\Omega(C(Y))$. If the functions in $\ker\varphi$ had no common
zero, compactness would give a finite subcover of the open sets
$\{y\in Y:f(y)\neq0\}$, with $f\in\ker\varphi$. There would then be
$f_1,\ldots,f_m\in\ker\varphi$ with no common zero. The corollary on the
involution implies that $f_j^*\in\ker\varphi$; hence
$g=\sum_{j=1}^m f_j^*f_j\in\ker\varphi$. But $g(y)>0$ for every $y\in Y$,
so $g$ is invertible in $C(Y)$ with inverse $1/g$. This contradicts
$1=\varphi(1)=\varphi(g)\varphi(g^{-1})=0$.
Thus there is $y_0\in Y$ at which every $f\in\ker\varphi$ vanishes.
For $h\in C(Y)$, the function $h-\varphi(h)1$ belongs to the kernel, so
$h(y_0)=\varphi(h)$. Thus $\varphi=\operatorname{ev}_{y_0}$.
The map $y\mapsto\operatorname{ev}_y$ is continuous because, for each
$h\in C(Y)$, its coordinate $y\mapsto\operatorname{ev}_y(h)=h(y)$ is continuous.
It is injective because $C(Y)$ separates points and surjective by the
preceding argument. Since $Y$ is compact and $\Omega(C(Y))$ is Hausdorff,
it is a homeomorphism.
If $F:\A\to C(Y)$ is a $*$-isomorphism, composition
$\chi\mapsto\chi\circ F$ is a homeomorphism from $\Omega(C(Y))$ to
$\Omega(\A)$, with inverse induced by $F^{-1}$. Therefore
$Y\cong\Omega(\A)=X$. Problem 2.5 develops this identification separately.
\end{proof}

\proofdiagram
{$\A$ commutative unital $C^*$; Gelfand--Naimark}
{$X=\Omega(\A)$}
{$\A\cong C(\Omega(\A))$; the characters of $C(X)$ are evaluations}
{The representation exists and $X$ is unique up to homeomorphism}

\begin{cthejemplo}{(Reconstructing the interval from its algebra)}
Consider $\A=C([0,1])$. The proof of Corollary~2.21 shows that every
character of $\A$ is an evaluation $\operatorname{ev}_x$ at a unique
$x\in[0,1]$. Thus the map $x\mapsto\operatorname{ev}_x$ identifies
$[0,1]$ topologically with $\Omega(\A)$.
If $\A$ were also $*$-isomorphic to $C(Y)$ for a compact Hausdorff space $Y$,
the same corollary would give $Y\cong\Omega(\A)\cong[0,1]$.
The function algebra therefore recovers the space from which it arises.
\end{cthejemplo}

\begin{investigacion}[unbreakable]
\textbf{Use in a construction.} In \cite[Sec.~2, pp.~465--466]{Karlovich2017},
$\SO(\mathbb R_+)$ is a unital commutative $C^*$-subalgebra of
$L^\infty(\mathbb R_+)$: its functions are continuous and bounded, and their
oscillation on $[r,2r]$ tends to zero as $r\to0$ or $r\to\infty$.
Identifying $C(\overline{\mathbb R}_+)$ with its restrictions to
$\mathbb R_+=(0,\infty)$, where $\overline{\mathbb R}_+=[0,\infty]$ is the
compactification, the fibers over $s\in\{0,\infty\}$ are
\[
 M_s(\SO(\mathbb R_+))=
 \{\xi\in\Omega(\SO(\mathbb R_+)):
 \xi|_{C(\overline{\mathbb R}_+)}=\operatorname{ev}_s\}.
\]
Proposition~2.2 of that source relates evaluation of a character in a fiber
to partial limits of slowly oscillating functions.
\textbf{What to recognize:} the characters and maximal ideal space studied
in the course serve here to encode information at the endpoints.
\end{investigacion}

\section{Continuous functional calculus}

\begin{cthteorema}{(Continuous functional calculus for a normal element)}
Let $\A$ be a unital $C^*$-algebra and $a\in\A$ normal. There is a unique
isometric $*$-homomorphism
\[
 C(\spec_\A(a))\longrightarrow\A,
 \qquad f\longmapsto f(a),
\]
that sends the identity function $z\mapsto z$ to $a$ and the constant
function $1$ to $1_\A$. Its image is $C^*(a,1_\A)$.
\end{cthteorema}
\origen{Theorem 2.22}

\begin{proof}[Proof]
Let $\B=C^*(a,1_\A)$. By Gelfand--Naimark,
$\Gamma_\B:\B\to C(\Omega(\B))$ is an isometric $*$-isomorphism.
The preceding proposition gives the homeomorphism
$\Psi:\Omega(\B)\to\spec(a)$, $\Psi(\varphi)=\varphi(a)$. The composition
\[
 C(\spec(a))\xrightarrow{f\mapsto f\circ\Psi} C(\Omega(\B))
 \xrightarrow{\Gamma_\B^{-1}}\B\subset\A
\]
is an isometric $*$-isomorphism from $C(\spec(a))$ onto $\B$.
For the identity function $\mathrm{id}(z)=z$, we have
$(\mathrm{id}\circ\Psi)(\varphi)=\varphi(a)=\widehat a(\varphi)$;
thus its image is $a$. The constant $1$ corresponds to
$\widehat{1_\A}$, so its image is $1_\A$.
If $F:C(\spec(a))\to\A$ is another isometric $*$-homomorphism with these
properties, then $F(\overline{\mathrm{id}})=a^*$ and it agrees with the
preceding construction on all polynomials in $z$ and $\overline z$.
Stone--Weierstrass makes that unital $*$-subalgebra dense in
$C(\spec(a))$; continuity of both maps extends equality to every
function. This gives the stated uniqueness and image.
\end{proof}

\proofdiagrambranch
{$\A$ is unital $C^*$; $a\in\A$ normal}
{$C^*(a,1_\A)$ commutative + Gelfand--Naimark}
{Identify $\Omega(C^*(a,1_\A))$ with $\spec_\A(a)$}
{Transport functions through the homeomorphism and the Gelfand transform}
{Define $f(a)$ uniquely}

\begin{cthejemplo}{(Spectral square root of a diagonal matrix)}
Take the positive normal matrix
$A=\operatorname{diag}(1,4)$ and the function
$f(t)=\sqrt t$ on $\spec(A)=\{1,4\}$. Theorem~2.22 gives
\[
 f(A)=\operatorname{diag}(f(1),f(4))
     =\operatorname{diag}(1,2),
 \qquad f(A)^2=A.
\]
In general, if
$A=U\operatorname{diag}(\lambda_1,\ldots,\lambda_n)U^*$ is normal,
$f(A)=U\operatorname{diag}(f(\lambda_1),\ldots,f(\lambda_n))U^*$.
The example shows that abstract functional calculus preserves the familiar
spectral recipe for matrices.
\end{cthejemplo}

\begin{cthnota}{(Dependence on the problems)}
Normality makes $C^*(a,1_\A)$ commutative; spectral permanence identifies
its spectrum with that of $a$ in $\A$. Problem 2.4 combines this calculus
with the spectrum of continuous functions from Problem 1.10 to obtain
$\spec(f(a))=f(\spec(a))$.
\end{cthnota}


\begin{cthdefinicion}{(Functional calculus by restriction from an ambient set)}
Let $\A$ be a unital $C^*$-algebra, $a\in\A$ normal, and $E\subset\mathbb C$
locally compact with $\spec_\A(a)\subset E$. For $f\in C_b(E)$ define
\[
 f(a):=(f|_{\spec_\A(a)})(a).
\]
This gives a unital $*$-homomorphism $C_b(E)\to\A$. The same formula
defines $f(a)$ if $f:E\to\mathbb C$ is continuous and unbounded, because its
restriction to the compact set $\spec_\A(a)$ is bounded. Only that
restriction matters.
\end{cthdefinicion}
\origen{Definition 2.23}

\begin{cthobservacion}{(Polynomials in $a$ and $a^*$)}
Let $a$ be normal in a unital $C^*$-algebra $\A$. If
$f(z)=z^2-\overline z^{\,3}+1$ for $z\in\mathbb C$, then
\[
 f(a)=a^2-(a^*)^3+1_\A.
\]
If $f(z)=|z|^2=\overline z z$, then $f(a)=a^*a$.
\end{cthobservacion}
\origen{Observation 2.24}

\begin{cthobservacion}{(In $C(X)$, functional calculus is composition)}
If $X$ is compact Hausdorff, $g\in C(X)$, and $f\in C(g(X))$, then
\[
 f(g)=f\circ g.
\]
\end{cthobservacion}
\origen{Observation 2.25}

\begin{proof}[Proof]
Under the identification $\Omega(C(X))\cong X$, the transform of $g$ takes
the value $g(x)$ at $x\in X$, and $\spec(g)=g(X)$. The construction of the
functional calculus transports $f$ through this transform; at each $x$ it
gives $f(g)(x)=f(g(x))$. Two functions in $C(X)$ that agree at every point
are equal, so $f(g)=f\circ g$.
\end{proof}

\begin{cthejemplo}{(The square root of $t^2$ is $|t|$)}
On $X=[-1,1]$ take the positive function $g(t)=t^2$.
Since $\spec_{C(X)}(g)=g(X)=[0,1]$, the continuous function
$f(s)=\sqrt s$ can be applied to $g$ by Theorem~2.22.
Observation~2.25 identifies the result pointwise:
\[
 \sqrt g(t)=f(g(t))=\sqrt{t^2}=|t|.
\]
The positive square root of the element $g$ is $|t|$, not the function $t$,
which changes sign.
\end{cthejemplo}

\begin{cthlema}{(Naturality of functional calculus)}
Let $\A,\B$ be unital $C^*$-algebras, $\Phi:\A\to\B$ a unital $*$-homomorphism,
$a\in\A$ normal, and $f\in C(\spec_\A(a))$. Since
$\spec_\B(\Phi(a))\subset\spec_\A(a)$, we evaluate $f$ at $\Phi(a)$ by
restricting it to $\spec_\B(\Phi(a))$. Then
\[
 \Phi(f(a))=f(\Phi(a)).
\]
\end{cthlema}
\origen{Lemma 2.26}

\begin{proof}[Proof]
The map $f\mapsto\Phi(f(a))$ is the composition of the functional calculus
for $a$ with $\Phi$; it is a unital, contractive $*$-homomorphism by the
preceding contractivity theorem.
The map $f\mapsto f(\Phi(a))$ is restriction
$C(\spec_\A(a))\to C(\spec_\B(\Phi(a)))$ followed by the functional calculus
for $\Phi(a)$, which is normal because $\Phi$ preserves products and
involution. It too is a unital, contractive $*$-homomorphism. Both maps send
the identity function to $\Phi(a)$ and the constant $1$ to $1_\B$; because
they preserve the involution, they also send $\overline z$ to $\Phi(a)^*$.
By linearity and multiplicativity they agree on polynomials in $z$ and
$\overline z$. These are dense by Stone--Weierstrass, and contractivity
allows us to pass to the limit for every $f\in C(\spec_\A(a))$.
\end{proof}

\proofdiagram
{$\Phi:\A\to\B$ unital and $*$; $a\in\A$ normal}
{Spectral inclusion; functional calculus; contractivity}
{The maps $f\mapsto\Phi(f(a))$ and $f\mapsto f(\Phi(a))$ agree on the identity and constants}
{By uniqueness/density, they agree on all of $C(\spec(a))$}

\begin{cthejemplo}{(A square root carried to multiplication operators)}
Let $\A=C([0,1])$ and represent it on $H=L^2([0,1])$ by
$\pi(f)=M_f$, the multiplication operator by $f$.
For $a(t)=t^2$, the positive square root in $\A$ is
$\sqrt a(t)=t$. By Lemma~2.26,
\[
 \pi(\sqrt a)=M_t=\sqrt{M_{t^2}}=\sqrt{\pi(a)}.
\]
Indeed, $M_t$ is positive and $M_t^2=M_{t^2}$.
We obtain the same operator by first taking the square root of the function
or by taking it after representing the algebra.
\end{cthejemplo}

\begin{cthteorema}{(A unital $*$-monomorphism is isometric)}
Let $\A,\B$ be unital $C^*$-algebras. If $\Phi:\A\to\B$ is a unital,
injective $*$-homomorphism, then
\[
 \|\Phi(a)\|=\|a\|\qquad(a\in\A).
\]
\end{cthteorema}
\origen{Theorem 2.27}

\begin{proof}[Proof]
We already know that $\|\Phi(a)\|\le\|a\|$. Suppose the inequality is strict
for some $a\neq0$ and put $x=a^*a$. Then $x$ is self-adjoint and
$\|\Phi(x)\|=\|\Phi(a)\|^2<\|a\|^2=\|x\|$.
Since $x$ is self-adjoint and $\|x\|=r(x)$, one of the points $\pm\|x\|$
belongs to $\spec(x)\subset[-\|x\|,\|x\|]$.
Moreover, $\spec(\Phi(x))\subset[-\|\Phi(x)\|,\|\Phi(x)\|]$.
Write $c=\|\Phi(x)\|$ and $d=\|x\|$, with $0\le c<d$, and define
\[
 f(t)=\max\!\left\{0,\min\!\left\{1,\frac{|t|-c}{d-c}\right\}\right\}
 \qquad(t\in\mathbb R).
\]
This function is continuous, vanishes on $[-c,c]$, and equals $1$ at both
points $\pm d$. Thus $f(\Phi(x))=0$, whereas $f(x)\neq0$ because the functional
calculus is isometric and $f$ does not vanish on $\spec(x)$. By naturality,
$0=f(\Phi(x))=\Phi(f(x))$, contradicting injectivity.
\end{proof}

\proofdiagrambranch
{$\Phi:\A\to\B$ unital, $*$, and injective}
{$a\in\A$; $x=a^*a$ is self-adjoint}
{If $\|\Phi(x)\|<\|x\|$, separate the two spectra by a continuous function}
{Naturality: $0=f(\Phi(x))=\Phi(f(x))$; injectivity contradicts $f(x)\neq0$}
{$\|\Phi(a)\|=\|a\|$}

\begin{cthejemplo}{(The diagonal representation of $C(\{1,2\})$)}
Let $X=\{1,2\}$ and $H=\ell^2(X)\cong\mathbb C^2$.
The unital representation $\pi:C(X)\to\mathcal B(H)$ assigns to
$f$ the matrix $M_f=\operatorname{diag}(f(1),f(2))$.
It is faithful, since $M_f=0$ implies $f(1)=f(2)=0$, and
\[
 \|\pi(f)\|=\max\{|f(1)|,|f(2)|\}=\|f\|_\infty.
\]
This concrete calculation agrees with the general conclusion of
Theorem~2.27: a faithful representation preserves the norm.
\end{cthejemplo}

\section{Recognizing self-adjoint elements, unitaries, projections, and positive elements by their spectra}

\begin{cthproposicion}{(Spectral classification of normal elements)}
Let $\A$ be a unital $C^*$-algebra and $a\in\A$ normal. Then
\begin{enumerate}[label=(\roman*)]
\item $a=a^*$ if and only if $\spec(a)\subset\mathbb R$;
\item $a$ is unitary if and only if $\spec(a)\subset\mathbb T$;
\item $a$ is a projection if and only if $\spec(a)\subset\{0,1\}$.
\end{enumerate}
\end{cthproposicion}
\origen{Proposition 2.28; corresponds directly to Problem 2.6}

\begin{proof}[Proof]
If $a=a^*$, we already know that $\spec(a)\subset\mathbb R$. Conversely,
if the spectrum is real, the functions $z$ and $\overline z$ agree on
$\spec(a)$, so functional calculus gives $a=a^*$.

If $a$ is unitary, the identity $a^*a=1_\A$ translates, by the isometric
functional calculus, into $|z|^2=1$ on $\spec(a)$; hence
$\spec(a)\subset\mathbb T$. Conversely, if
$\spec(a)\subset\mathbb T$, then $|z|^2=1$ on the spectrum and
$a^*a=(\overline z z)(a)=1_\A$; normality also gives $aa^*=1_\A$.

If $a$ is a projection, $a^2-a=0$. The isometry of functional calculus
applied to $z\mapsto z^2-z$ gives $z^2-z=0$ throughout $\spec(a)$;
therefore $\spec(a)\subset\{0,1\}$. Conversely, if the spectrum is
contained in $\{0,1\}$, then $z^2-z=0$ and $\overline z-z=0$ on the
spectrum; functional calculus gives $a^2=a$ and $a^*=a$.
\end{proof}

\proofdiagrambranch
{$\A$ is unital $C^*$; $a\in\A$ normal}
{Continuous functional calculus}
{Translate $a^*=a$, $a^*a=1$, $a^2=a$ into identities of functions on $\spec(a)$}
{$\overline z=z$, $|z|^2=1$, $z^2=z$ on the spectrum}
{The spectral inclusions characterize the three classes}

\begin{cthejemplo}{(Counterexamples: a small spectrum without normality)}
In $M_2(\mathbb C)$, take
$N=\begin{pmatrix}0&1\\0&0\end{pmatrix}$ and
$J=\begin{pmatrix}1&1\\0&1\end{pmatrix}$.
Both are triangular, so
$\spec(N)=\{0\}\subset\mathbb R\cap\{0,1\}$ and
$\spec(J)=\{1\}\subset\mathbb T$.
However, $N^*\neq N$ and $N^2=0\neq N$;
moreover, $J^*J=\begin{pmatrix}1&1\\1&2\end{pmatrix}\neq I$.
Problem~2.6(d) develops these cases. They show that the converse conclusions
of Proposition~2.28 require normality.
\end{cthejemplo}

\begin{cthproposicion}{(Four unitaries linearly generate every element)}
Let $\A$ be a unital $C^*$-algebra. Every $a\in\A$ is a linear combination
of at most four unitary elements.
\end{cthproposicion}
\origen{Proposition 2.29}

\begin{proof}[Proof]
Write $a=x_1+ix_2$ with $x_j$ self-adjoint. Let
$r_j=\max\{1,\|x_j\|\}$ and $h_j=x_j/r_j$, so $\|h_j\|\le1$.
Since $h_j=h_j^*$, its spectrum lies in $[-1,1]$. The continuous function
$f(t)=t+i\sqrt{1-t^2}$ satisfies
$|f(t)|^2=t^2+(1-t^2)=1$ and $f(t)+\overline{f(t)}=2t$ on that interval.
By functional calculus, $u_j=f(h_j)$ satisfies
$u_j^*u_j=u_ju_j^*=(|f|^2)(h_j)=1_\A$ and is unitary. Moreover,
$u_j+u_j^*=(f+\overline f)(h_j)=2h_j$. Therefore
\[
 a=\frac{r_1}{2}(u_1+u_1^*)+\frac{i r_2}{2}(u_2+u_2^*),
\]
the asserted combination.
\end{proof}

\proofdiagram
{$\A$ is unital $C^*$; $a=x_1+ix_2$ with $x_j=x_j^*$}
{Scale $h_j=x_j/r_j$ with $r_j=\max\{1,\|x_j\|\}$}
{$u_j=f(h_j)$, $f(t)=t+i\sqrt{1-t^2}$; $x_j=r_j(u_j+u_j^*)/2$}
{$a$ is a combination of $u_1,u_1^*,u_2,u_2^*$}

\begin{cthejemplo}{(A self-adjoint element as the average of two unitaries)}
In $M_2(\mathbb C)$ consider
$H=\operatorname{diag}(1/2,-1/2)$.
It is self-adjoint and $\|H\|\le1$. The construction in
Proposition~2.29 produces
\[
 U=H+i\sqrt{I-H^2}
 =\operatorname{diag}\!\left(
 \frac12+i\frac{\sqrt3}{2},
 -\frac12+i\frac{\sqrt3}{2}\right).
\]
Each diagonal entry has modulus one, so $U$ is unitary, and
$H=(U+U^*)/2$. This case exhibits the step that allows us to express
any element as a combination of at most four unitaries.
\end{cthejemplo}

\begin{cthproposicion}{(Positive square root)}
Let $\A$ be a unital $C^*$-algebra and $a\in\A$ normal with
$\spec_\A(a)\subset[0,\infty)$. Then there exists a positive self-adjoint
element $y\in C^*(a,1_\A)$ such that
\[
 y^2=a.
\]
The root obtained by functional calculus is denoted $y=\sqrt a$.
\end{cthproposicion}
\origen{Proposition 2.30}

\begin{proof}[Proof]
The function $t\mapsto\sqrt t$ is continuous and real on $\spec(a)$.
Define $y=\sqrt{\,\cdot\,}(a)$. Then $y^*=y$ and
$y^2=(t\mapsto t)(a)=a$.
The identification $C^*(a,1_\A)\cong C(\spec_\A(a))$ shows that the
transform of $y$ is $t\mapsto\sqrt t\ge0$; by the description of the
spectrum in a function algebra and spectral permanence,
$\spec_\A(y)\subset[0,\infty)$.
\end{proof}

\proofdiagram
{$\A$ is unital $C^*$; $a\in\A$ normal and $\spec_\A(a)\subset[0,\infty)$}
{Continuous functional calculus}
{Apply $f(t)=\sqrt t$}
{$y=f(a)$ is self-adjoint, $\spec(y)\subset[0,\infty)$, and $y^2=a$}

\begin{cthejemplo}{(Modulus of a nilpotent matrix)}
Let $A=\begin{pmatrix}0&2\\0&0\end{pmatrix}$ in $M_2(\mathbb C)$.
Although $A^2=0$, the product
\[
 A^*A=\operatorname{diag}(0,4)
\]
is positive. Proposition~2.30 allows us to take its positive square root,
which is calculated directly here:
\[
 |A|=(A^*A)^{1/2}=\operatorname{diag}(0,2).
\]
This $|A|$ retains the magnitude of the action of $A$ and is the object used
in Problem~2.7 for the polar decomposition. In general, applying the
proposition to $a^*a$ requires the standard fact $a^*a\ge0$.
\end{cthejemplo}

\begin{investigacion}
\textbf{Direct example.} In \cite[Sec.~1, p.~20]{BFK2006}, the operators
$(U_z\varphi)(t)=\varphi(zt)$, with $z,t\in\mathbb T$ and
$\varphi\in L^2(\mathbb T)$, form a unitary representation of the group
$\mathbb T$ on $L^2(\mathbb T)$; rotations preserve arc-length measure.
\textbf{What to recognize:} the abstract unitary here implements a concrete
geometric symmetry and participates in the nonlocal extension of the algebra
studied in that source.
\end{investigacion}

\section*{Historical note}
\origen{Historical note of the source chapter, pp.~8--10}

\paragraph{The original axioms.}
Gelfand and Naimark's 1943 work sought to represent $C^*$-algebras,
including noncommutative ones, by bounded operators on Hilbert spaces.
The source notes recall that the original definition included, besides
$\|a^*a\|=\|a\|^2$, isometry of the involution and invertibility of
$1+a^*a$. It was later proved that these additional conditions follow from
the $C^*$-identity.

\paragraph{$B^*$ and $C^*$.}
For a time, $B^*$ denoted the axiomatic structure and $C^*$ denoted
adjoint-closed subalgebras of bounded operators on Hilbert space. The
general Gelfand--Naimark representation theorem identifies the two
descriptions; the commutative theorem proved here is the course's step
toward that idea.

\paragraph{Rigidity and functional calculus.}
Theorem~2.6, Corollary~2.7, and Theorem~2.27 in the source notes show three
forms of norm rigidity. Spectral permanence contrasts with the disk algebra
of Problem~1.12, which is a closed subalgebra of $C(\mathbb T)$ but not a
$C^*$-subalgebra. Functional calculus transports continuous functions on
the spectrum to elements of the algebra without first choosing a
representation on Hilbert space; Stone--Weierstrass explains why functions
of $a$ and $a^*$ suffice.

\paragraph{Topology and positivity.}
The commutative classification recovers the compact space $X$ from the
characters of $C(X)$. This viewpoint motivated the later language of
noncommutative topology. The square root in Proposition~2.30 of the source
notes opens the systematic study of positive elements and order, while
Proposition~2.29 there anticipates the importance of unitaries for the
geometry of the algebra.
